\subsection*{\texttt{\_\_co\_\_ auto} 返回类型}

若结果形状更适合在函数体中描述，可以使用 \co@auto@。对于 \co@__co__ auto matmul(...)@，编译器从 \co@return output@ 所返回的缓冲区推导类型。返回值仍然有元素类型和形状，\co@auto@ 只是省去了签名中重复的声明。若希望直接在接口中约定输出形状，则可以写出完整返回类型。

逐元素算术、标量矩阵乘法以及返回缓冲区的 MMA 核函数，都遵循相同的返回约定。函数内部怎样计算元素，并不影响这一规则。
